% 基礎となる斉次Sklyanin代数からPBW・CY入力を導く。
\section{定数項変形のPBW性と自己双対な三次元解消}
\label{sec:quotient-pbw-cy}
以下では斉次Sklyanin代数
\[
 \begin{gathered}
 S_1=\CC\langle X,Y,Z\rangle/(r_X,r_Y,r_Z),\\
 r_X=YZ-ZY-X^2,\quad r_Y=ZX-XZ-Y^2,\quad r_Z=XY-YX-Z^2
 \end{gathered}
\]
のKoszul AS正則性（次元3、Hilbert級数$(1-t)^{-3}$）を既知入力とする。
この斉次代数の基本的な性質は\cite[Proposition~2.21, \S8.2]{WWYSklyanin}を参照する。
非斉次代数のPBW性とCY性を独立した未証明の仮定として追加する必要はない。
これらは\cite[Theorem~3.3.2]{EGDelPezzo}の既知の一般的なpotential変形定理の特殊化でもある。
以下では具体的なsyzygyと自己双対解消による証明、および被約基底上での相対化を与える。

\begin{proposition}[任意の定数項変形のPBW性]
\label{prop:quotient-constant-pbw}
任意の$c_X,c_Y,c_Z\in\CC$に対し
\[
 R(c)=\CC\langle X,Y,Z\rangle/
 (r_X-c_X,r_Y-c_Y,r_Z-c_Z)
\]
は語長filtrationについて$\operatorname{gr}R(c)=S_1$を満たす。
特に$R_1=R(1,0,0)$について$\dim F_nR_1=\binom{n+3}{3}$である。
\end{proposition}
\begin{proof}
$V=\langle X,Y,Z\rangle$、$\mathcal R=\langle r_X,r_Y,r_Z\rangle$とする。
Koszul解消から$(V\otimes\mathcal R)\cap(\mathcal R\otimes V)$は一次元で、
\[
 \eta=X\otimes r_X+Y\otimes r_Y+Z\otimes r_Z
     =r_X\otimes X+r_Y\otimes Y+r_Z\otimes Z
\]
で生成される。最後の等式は六つの三文字語と三つの立方の展開から直接従う。
Braverman--Gaitsgory\cite[Main Theorem~0.5, Lemma~3.3]{BGPBW}の記号で
一次補正は零、定数補正は$\beta(r_i)=c_i$である。
同定理の条件(i),(iii)は自動的に成立し、条件(ii)も
\[
 (\beta\otimes1-1\otimes\beta)(\eta)
 =c_XX+c_YY+c_ZZ-Xc_X-Yc_Y-Zc_Z=0
\]
から従う。従ってPBW性が成立する。
\end{proof}

\begin{proposition}[被約有限型基底上の相対PBW性]
\label{prop:quotient-relative-pbw}
$C$を被約有限生成可換$\CC$-代数、$c_X,c_Y,c_Z\in C$とする。
同じ表示で定義した$C$-代数$R_C(c)$は
\[
 \operatorname{gr}R_C(c)\simeq C\otimes_\CC S_1
\]
を満たし、斉次語基底を持ち上げると$C$上自由である。
\end{proposition}
\begin{proof}
$V_n$を長さ$n$以下の自由語の$C$加群に、
$|a|+|b|\le n-2$の関係式倍$a(r_i-c_i)b$を課した有限表示加群とする。
各閉点への特殊化は前命題のPBW代数であり、$\dim(V_n\otimes k(s))=\binom{n+3}{3}$。
ここでPBW性から、次数$n$以下の関係はこれら有界次数の倍で生成される。
実際、関係の最高次項を斉次関係の倍で消して次数を下げればよい。

$V_n$の有限表示行列のrankは全閉点で一定となる。
それより大きいminorは被約性とNullstellensatzで$C$上零であり、
最大rankのminorの非消滅開集合は$\operatorname{Spec}C$を覆う。
各開集合で行列を恒等blockと零blockに変形すると、$V_n$は有限射影となり、
形成は任意の基底変更と可換する。
$V_n\to V_{n+1}$は全閉点で単射で、rankも一定なので、同じ議論から
単射で有限射影なcokernelを持つ。

$\varinjlim V_n=R_C(c)$なので$V_n$はその実際の語長filtrationに等しい。
自然な全射$C\otimes(S_1)_n\to V_n/V_{n-1}$は
全閉点で同型であり、有限射影加群の同rankの写像なので同型である。
各商は自由だからfiltrationの短完全列は分裂し、語基底の持上げが
$C$上の自由基底を与える。
被約性を外して閉点のrankだけから同じ結論を推論してはいけない。
\end{proof}

\begin{proposition}[定数項変形の三次元Calabi--Yau性]
\label{prop:quotient-constant-cy}
$R(c)$はhomologically smoothで、
\[
 \mathbf R\operatorname{Hom}_{R(c)^e}(R(c),R(c)^e)
 \simeq R(c)[-3].
\]
従って特に$R_1$は三次元Calabi--Yauである。
\end{proposition}
\begin{proof}
$A=R(c)$とし、$P_i=A\otimes E_i\otimes A$を、
$E_0=\CC$、$E_1=V$、$E_2=\mathcal R$、$E_3=\CC\eta$で定める。
微分を
\[
 d_1(1\otimes X_i\otimes1)=X_i\otimes1-1\otimes X_i,
\]
$d_2(1\otimes r_i\otimes1)=\partial(r_i)$、
\[
 d_3(1\otimes\eta\otimes1)
 =\sum_i(X_i\otimes r_i\otimes1-1\otimes r_i\otimes X_i)
\]
とする。ここで$\partial(w)$は、語$w$の各文字を順に真ん中の$V$へ移し、
その前後の語を両端へ置く普遍微分である。
$d_1d_2=0$は$\partial(r_i)$のtelescopingと$r_i=c_i$から、
$d_2d_3=0$は$\eta$の二表示と定数の中心性から従う。
各$E_i$に次数$i$を与えたfiltrationのassociated gradedは
$S_1$のKoszul双加群解消である。
最高次項を消して次数を下げる帰納法により
\[
 0\to P_3\to P_2\to P_1\to P_0\to A\to0
\]
は完全であり、有限自由双加群解消となる。

自己双対性も直接確認できる。
$L_i=X_i\otimes1$、$R_i=1\otimes X_i^{\mathrm{op}}$、
$D_i=L_i-R_i$と書くと、$d_2$の各基底像の係数行列は
\[
 \begin{pmatrix}
 -(L_X+R_X)&D_Z&-D_Y\\
 -D_Z&-(L_Y+R_Y)&D_X\\
 D_Y&-D_X&-(L_Z+R_Z)
 \end{pmatrix}.
\]
この行列は、両端を交換する反同型$\tau(L_i)=R_i$と転置の合成で不変である。
$d_1$と$d_3$も同じ交換で符号を除き互いに双対となる。
従って上の複体の$A^e$-双対は、標準の複体の符号を調整して
元の複体を三つずらしたものとなる。
右端の余核の双加群は通常の$A$であり、Nakayama twistは生じない。
これで表示した導来双対とsmoothnessを同時に得る。
\end{proof}

\begin{corollary}[指定混合族の周囲と中心三次商]
\label{cor:quotient-ambient-pbw}
全ての$\alpha\in\CC$について
\[
 B_\alpha=\CC[u,v]\langle x,y,z\rangle/
 ([x,y]-z^2,[y,z]-x^2-\alpha uv,[z,x]-y^2)
\]
はnoetherianなKoszul AS正則整域（次元5、index 5）である。
$\Omega_1$を本文の対称化中心三次式とすると
$W_\alpha=\Omega_1+\alpha uvx+u^3+v^3$は中心正則元で、
$A_\alpha=B_\alpha/(W_\alpha)$は$\CC[u,v]$上自由であり、
全ての非零一次式$au+bv$は$A_\alpha$上正則である。
また$B_\alpha/(v-\lambda u)$は全ての$\lambda\in\CC$で
Koszul AS正則（次元4、index 4）である。
\end{corollary}
\begin{proof}
相対PBW命題を$C=\CC[u,v]$、$c_X=\alpha uv,c_Y=c_Z=0$に適用する。
$u,v$は中心正則列で$B_\alpha/(u,v)=S_1$。
相対filtrationのassociated gradedが$S_1[u,v]$なので
整域性とnoetherian性も従う。

次数1の中心正則元$z$に対するchange-of-rings Tor spectral sequenceを使う。
$B/(z)$がKoszulなら、$\operatorname{Tor}^B_q(B/(z),\CC)$は
$q=0,1$で内部次数$q$に集中するので、$E_2$は全て内部次数$p+q$に集中する。
従って$B$もKoszulで、$\operatorname{pd}_B\CC=
\operatorname{pd}_{B/(z)}\CC+1$となる。
これを$u,v$に適用して両側の射影次元5を得る。
minimal graded解消とgood filtrationで全加群の大域次元も5となる。
中心正則列に対するReesの補題からAS--Gorenstein条件とindex 5も従う。
同じ論法を基底$\CC[u]$、$c_X=\alpha\lambda u^2$に適用すれば
$B_\alpha/(v-\lambda u)$について次元4、index 4を得る。

$\Omega_1+mx$の中心性は$m=1$での厳密な中心元計算
（命題\ref{prop:quotient-verified-center6}）と、
$m\ne0$での$(x,y,z)=\sqrt m(X,Y,Z)$による規格化から従う。
$m=0$では斉次中心元である。
相対PBW基底の各係数について全閉点でこの恒等式が成立するので、
被約基底$\CC[u,v]$上でも中心的である。
$W_\alpha$は非零で、最高$x,y,z$次数部分$\Omega_1$は$S_1$上正則。
従って
\[
 \operatorname{gr}_{x,y,z}A_\alpha
 =(S_1/(\Omega_1))\otimes\CC[u,v].
\]
斉次語基底の持上げから最後の自由性と正則性を得る。
\end{proof}

\begin{proposition}[指定混合族の全係数値でのNakayama twistの消滅]
\label{prop:quotient-ambient-nakayama}
全ての$\alpha\in\CC$について$B_\alpha$のNakayama自己同型は恒等であり、
$B_\alpha$は五次元Calabi--Yau代数である。
$A_\alpha$と$C_\alpha=A_\alpha/(v-u)$のbalanced dualizing complexは
それぞれ$A_\alpha(-2)[4]$、$C_\alpha(-1)[3]$である。
従って球面twistの構成に必要なNakayama twistの恒等性も追加仮定ではない。
\end{proposition}
\begin{proof}
$E=B_\alpha^!$をKoszul双対とし、生成元を$e_x,e_y,e_z,e_u,e_v$とする。
直交関係を取ると、15本の二次関係は
\[
 e_x^2+e_ye_z=e_y^2+e_ze_x=e_z^2+e_xe_y=0,
 \qquad e_u^2=e_v^2=0,
\]
\[
 e_ie_j+e_je_i=0\quad(i<j,\ (i,j)\ne(u,v)),\qquad
 e_ue_v+e_ve_u+\alpha e_ye_z=0.
\]
$E$はHilbert級数$(1+t)^5$のgraded Frobenius代数である。
以下の有限計算は$\mathbb Q[\alpha]$上でその構造を直接検証する。
次数4の基底を
\[
 b_1=e_x^3e_u,\ b_2=e_x^3e_v,\ b_3=e_x^2e_ue_v,\
 b_4=e_xe_ye_ue_v,\ b_5=e_xe_ze_ue_v
\]
とし、次数5の基底を$h=e_x^3e_ue_v$とする。
$\lambda:E_5\to\CC$を$\lambda(h)=1$で定めると、
両方の行列$(\lambda(e_i b_j))_{i,j}$と
$(\lambda(b_j e_i))_{i,j}$は同じ
\[
 M=\begin{pmatrix}
 0&0&1&0&0\\
 0&0&0&0&1\\
 0&0&0&-1&0\\
 0&-1&0&0&0\\
 1&0&0&0&0
 \end{pmatrix},\qquad\det M=1
\]
となる。従ってFrobeniusのNakayama自己同型$\nu_E$は次数1で恒等、
次数1が代数を生成するため全体で恒等である。

再現コード\texttt{computations/verify\_central\_quotient\_koszul.py}は
次数6までの自由語商を構成し、次元$1,5,10,10,5,1,0$、
全ての相補次数のFrobenius対合の行列式$1,1,-1,-1,1,1$、
上の左右の積の一致を厳密に検査する。
全ての行簡約pivotは非零有理定数であり、$\alpha$を含む式で割らない。
従って基底と等式は全ての複素$\alpha$へ特殊化でき、一般点だけの計算ではない。
同コードは三次syzygyと前命題の自己双対な双加群複体の積も検査する。

Koszul AS正則代数の双対公式
\cite[proof of Theorem~6.3]{RRZNakayama}は
$\mu_{B_\alpha}|_{B_1}=(-1)^{5+1}(\nu_E|_{E_1})^*$を与える。
符号も正なので$\mu_{B_\alpha}=\mathrm{id}$。
\cite[Lemma~1.2]{RRZNakayama}よりAS正則性はskew CY性を与えるから、
これは通常の五次元CY性である。
中心正則元$W_\alpha$と$v-u$による商には
\cite[Lemma~1.5]{RRZNakayama}を適用でき、normalizing automorphismも恒等なので
両商のNakayama自己同型も恒等となる。
AS次元とindexは順に$(4,2)$、$(3,1)$であり、
\cite[Corollary~4.3 and its preceding discussion]{YZSerre}の
balanced dualizing complexの式を適用して結論する。
\end{proof}
